\section{Sudoku Solver}
\label{sec:rec-sudoku-solver}

\problemstatement{Given a partially filled $9\times9$ Sudoku board
(empty cells marked \texttt{'.'}), fill every empty cell with a digit
from $1$ to $9$ so that every row, every column, and every one of the
nine $3\times3$ boxes contains each digit exactly once. It is guaranteed
the input has a unique solution.}

\begin{patternbox}
This closing backtracking problem combines two ideas already
established: N-Queens's (Section~\ref{sec:rec-n-queens}) ``track
constraint membership in auxiliary structures for fast checks''
(generalized here to three overlapping constraints -- row, column, and
box -- instead of just column and two diagonals), and a new control-flow
idea needed because this problem wants \textbf{any one} valid completion,
not \emph{every} one: the recursion must \textbf{return a boolean}
signaling success, so that the instant a complete valid filling is found
anywhere in the search, every enclosing call can immediately stop trying
further alternatives and propagate that success straight back to the
top.
\end{patternbox}

\subsection*{Understanding the problem carefully}

Scan the board for the next empty cell (in, say, row-major order). For
that cell, try each digit $1$ through $9$; a digit is valid if it does
not already appear elsewhere in the same row, the same column, or the
same $3\times3$ box. If a valid digit is found, place it and recursively
attempt to solve the rest of the board. If that recursive attempt
\emph{succeeds} (returns \texttt{true}, meaning a full valid solution was
found somewhere beneath this choice), we are done -- stop trying further
digits at this cell and propagate \texttt{true} upward immediately. If
it \emph{fails} (returns \texttt{false}), undo this digit and try the
next one. If no digit from $1$ to $9$ works at this cell, this entire
branch is a dead end -- return \texttt{false}, triggering the caller to
backtrack and try a different digit at an earlier cell.

\begin{ideabox}[Boolean-Returning Backtracking: Stop on First Success]
$\textsc{Solve}(\textit{board})$: find the next empty cell
$(\textit{row},\textit{col})$; if none exists, the board is completely
and validly filled -- return \texttt{true} (base case, success).
Otherwise, for $d$ from $1$ to $9$: if placing $d$ at
$(\textit{row},\textit{col})$ is valid: place it; if
$\textsc{Solve}(\textit{board})$ returns \texttt{true}: return
\texttt{true} immediately (propagate success, no undo needed -- this
digit is part of the final solution). Otherwise: remove $d$ (undo) and
try the next digit. If every digit from $1$ to $9$ has been tried
without success: return \texttt{false}.
\end{ideabox}

\subsection*{Why a boolean return value, with immediate propagation, is the key new idea}

Section~\ref{sec:rec-n-queens} and earlier sections collected \emph{all}
valid solutions, so every branch had to be explored regardless of
earlier successes. Sudoku, by contrast, only needs \emph{one} valid
completion (and the problem guarantees exactly one exists), so the
moment a recursive call confirms a digit choice leads all the way to a
complete, valid board, there is no reason to keep searching -- every
enclosing call should immediately stop and also report success,
unwinding the entire call stack without exploring any further
alternatives. This is exactly Section~\ref{sec:rec-check-subsequence-sum-k}'s
short-circuit-on-success idea, now applied to a problem where the
``success'' condition is ``the rest of the board can be completed,''
checked recursively rather than as a simple base-case comparison.

\subsection*{Algorithm}

\begin{enumerate}
  \item $\textsc{Solve}(\textit{board})$:
  \begin{enumerate}
    \item Find the next empty cell $(\textit{row}, \textit{col})$; if
          none, return \texttt{true}.
    \item For $d$ from $1$ to $9$:
    \begin{itemize}
      \item If $d$ is valid at $(\textit{row}, \textit{col})$ (no
            conflict in its row, column, or $3\times3$ box): place
            $d$. If $\textsc{Solve}(\textit{board})$ returns
            \texttt{true}: return \texttt{true}. Else: remove $d$
            (undo).
    \end{itemize}
    \item Return \texttt{false} (no digit worked here; trigger
          backtracking).
  \end{enumerate}
\end{enumerate}

\subsection*{Dry run}

\dryrun{A small conceptual trace: suppose cell $(0,2)$ is empty, and
digits $1$--$5$ are already ruled out by its row/column/box, but
$6$ is valid. We place $6$ and recurse to the next empty cell. Several
cells later, suppose a dead end is reached where no digit from $1$--$9$
is valid at some cell $(4,7)$ -- this returns \texttt{false}, which
propagates back to the most recent choice point (perhaps the placement
at $(3,5)$), forcing it to try its \emph{next} candidate digit instead
of the one that led to this dead end. This unwind-and-retry continues
until either a different early choice avoids the dead end entirely (and
the search eventually completes successfully, propagating \texttt{true}
all the way back to the top), or -- if the puzzle is unsolvable -- every
possibility is exhausted and the top-level call returns \texttt{false}.}

\subsection*{C++ implementation}

\begin{lstlisting}
#include <bits/stdc++.h>
using namespace std;

bool isValid(vector<vector<char>>& board, int row, int col, char d) {
    for (int i = 0; i < 9; i++) {
        if (board[row][i] == d) return false;
        if (board[i][col] == d) return false;
        if (board[3*(row/3) + i/3][3*(col/3) + i%3] == d) return false;
    }
    return true;
}

bool solve(vector<vector<char>>& board) {
    for (int row = 0; row < 9; row++) {
        for (int col = 0; col < 9; col++) {
            if (board[row][col] != '.') continue;

            for (char d = '1'; d <= '9'; d++) {
                if (isValid(board, row, col, d)) {
                    board[row][col] = d;
                    if (solve(board)) return true;
                    board[row][col] = '.';
                }
            }
            return false; // no digit worked at this empty cell
        }
    }
    return true; // no empty cell left: solved
}

int main() {
    vector<vector<char>> board = {
        {'5','3','.','.','7','.','.','.','.'},
        {'6','.','.','1','9','5','.','.','.'},
        {'.','9','8','.','.','.','.','6','.'},
        {'8','.','.','.','6','.','.','.','3'},
        {'4','.','.','8','.','3','.','.','1'},
        {'7','.','.','.','2','.','.','.','6'},
        {'.','6','.','.','.','.','2','8','.'},
        {'.','.','.','4','1','9','.','.','5'},
        {'.','.','.','.','8','.','.','7','9'}
    };

    solve(board);

    for (auto& row : board) {
        for (char c : row) cout << c << " ";
        cout << endl;
    }
    return 0;
}
\end{lstlisting}

\begin{complexitybox}
\textbf{Time Complexity.} Exponential in the worst case, roughly
$O(9^{m})$ where $m$ is the number of empty cells (each could in
principle try $9$ digits), though the row/column/box validity checks
prune the vast majority of branches in practice, making real Sudoku
puzzles solve very quickly.
\textbf{Space Complexity.} $O(m)$ for the recursion depth (at most one
call per empty cell along the current path), plus $O(1)$ extra for the
$\textit{isValid}$ check (which scans a fixed $9$ cells each for row,
column, and box).
\end{complexitybox}

\begin{takeawaybox}
Boolean-returning backtracking -- where finding \emph{one} success
anywhere in a subtree should immediately propagate \texttt{true} all the
way up, aborting every remaining sibling branch -- is the right shape
whenever a problem asks for any single valid solution (or asks whether
one exists at all), as opposed to enumerating every solution
(N-Queens) or counting them. Recognizing which of these three modes
(generate all, exists, count) a backtracking problem needs, exactly as
flagged back in Section~\ref{sec:rec-count-subsequence-sum-k}'s
takeaway, remains the first and most important design decision.
\end{takeawaybox}
